\section*{الفصل \ref{ch:b1:atmospheric-circulation} --- الدوران العام للغلاف الجوي}
% القيم: calc/ch28.py

\begin{solution}{exo:b1:atmospheric-circulation:1}
(أ)~$f = 2\Omega\sin\varphi$: $0$ عند خطّ الاستواء، و\qty{7.29e-5}{s^{-1}} عند \ang{30}،
\qty{1.031e-4}{s^{-1}} عند \ang{45}، و\qty{1.458e-4}{s^{-1}} عند القطب. (ب)~$fV = \num{1.031e-4} \times 10 = \qty{1.03e-3}{m.s^{-2}}$، نحو \num{1.05e-4}
من $g$. (ج)~في نصف الكرة الشمالي ينحرف الاتجاه إلى يمين الحركة: ريح تهبّ نحو الشمال تُدفع نحو الشرق.
\end{solution}

\begin{solution}{exo:b1:atmospheric-circulation:2}
(أ)~$|\nabla p|/\rho = (400/\num{4e5})/1.225 = \qty{8.16e-4}{m.s^{-2}}$.
(ب)~$10/\num{8.16e-4} = \qty{12250}{s}$، أي \qty{3.4}{h}. (ج)~تسارع كوريوليس لريح قدرها \qty{10}{m/s} عند \ang{45} شمالًا هو \qty{1.03e-3}{m.s^{-2}}: الاثنان بالحجم نفسه (النسبة $0.79$). وخلال ساعات قليلة من بدء الحركة يشعر الهواء بقوة كوريوليس قادرة على موازنة تدرّج الضغط: فتستقرّ الريح قرب السرعة الجيوستروفية، هنا $\num{8.16e-4}/\num{1.031e-4} = \qty{7.9}{m/s}$، بدلًا من أن تتسارع على طول التدرّج.
\end{solution}

\begin{solution}{exo:b1:atmospheric-circulation:3}
(أ)~$f = 2 \times \num{7.292e-5} \times \sin\ang{40} = \qty{9.374e-5}{s^{-1}}$؛
$|\vec v_g| = (400/\num{2.5e5})/(1.225 \times \num{9.374e-5}) = \qty{13.9}{m/s}$، أو
$13.93/0.5144 = \qty{27.1}{kn}$. (ب)~في نصف الكرة الشمالي تُبقي الريح الضغط المنخفض على يسارها: مع كون المنخفض جنوبًا، تهبّ الريح نحو الغرب، من الشرق (\omterm{def:b1:spherical-astronomy:horizontal}{السمت} \ang{90}). (ج)~تشير العصا شرقًا من دائرة المحطة، نحو الريح؛ و\qty{27.1}{kn} تُقرَّب إلى \qty{25}{kn}، ريشتان كاملتان ونصف ريشة، تُرسم على جانب العصا مع عقارب الساعة، هنا جانبها الجنوبي، نحو المنخفض.
\end{solution}

\begin{solution}{exo:b1:atmospheric-circulation:4}
(أ)~\ang{15} شمالًا: \omterm{def:b1:atmospheric-circulation:winds}{الرياح التجارية} الشمالية الشرقية؛ \ang{15} جنوبًا: \omterm{def:b1:atmospheric-circulation:winds}{الرياح التجارية} الجنوبية الشرقية؛
\ang{45} شمالًا و\ang{45} جنوبًا: \omterm{def:b1:atmospheric-circulation:winds}{الرياح الغربية}. (ب)~عند \ang{15} يعود هواء السطح \omterm{def:b1:atmospheric-circulation:cells}{لخليّة هادلي} نحو خطّ الاستواء، وتحرفه قوة كوريوليس نحو الغرب (يمينًا في الشمال، ويسارًا في الجنوب): شرقية، ب$u$ سالبة. وعند \ang{45} يتحرّك هواء السطح \omterm{def:b1:atmospheric-circulation:cells}{لخليّة فيريل} نحو القطب ويُحرف نحو الشرق: غربية، ب$u$ موجبة. (ج)~من الخطوط العرضية الثلاثة، يبلغ الضغط أعلاه عند \ang{30} شمالًا، \qty{1016.8}{hPa}: يقع \omterm{def:b1:atmospheric-circulation:cells}{المرتفع شبه المداري} قرب \ang{30} شمالًا (يبلغ الملف الكامل ذروته قرب \ang{35} شمالًا).
\end{solution}

\begin{solution}{exo:b1:atmospheric-circulation:5}
(أ)~$|\nabla p|/\rho = (200/\num{5e5})/1.225 = \qty{3.27e-4}{m.s^{-2}}$، مقسومًا على
$f$: \qty{2.6}{m/s} عند \ang{60}، و\qty{2.9}{m/s} عند \ang{50}، و\qty{4.5}{m/s} عند
\ang{30}، و\qty{6.5}{m/s} عند \ang{20}، و\qty{12.9}{m/s} عند \ang{10}، و
\qty{25.7}{m/s} عند \ang{5}. (ب)~تدور قوة كوريوليس الريح في زمن من رتبة
$1/f$: \qty{2.2}{h} عند \ang{60}، و\qty{11}{h} عند \ang{10}، و\qty{22}{h} عند \ang{5}، وطويل إلى ما لا نهاية عند خطّ الاستواء. وقرب خطّ الاستواء يتغيّر حقل الضغط، ويعمل الاحتكاك، قبل زمن طويل من أن تُدار الريح على طول الخطوط المتساوية الضغط: فالتوازن لا يحدث أبدًا، ويتدفّق الهواء جزئيًّا نحو أسفل التدرّج، بسرعات متواضعة.
\end{solution}

\begin{solution}{exo:b1:atmospheric-circulation:6}
(أ)~$u = 464.6\,\sin^2\varphi/\cos\varphi$: \qty{14.2}{m/s} عند \ang{10}،
\qty{32.2}{m/s} عند \ang{15}، و\qty{57.8}{m/s} عند \ang{20}، و\qty{134.1}{m/s} عند
\ang{30}. (ب)~$u = 464.6\,(\cos^2\ang{10} - \cos^2\varphi)/\cos\varphi$:
\qty{17.7}{m/s} عند \ang{15}، و\qty{42.9}{m/s} عند \ang{20}، و\qty{117.9}{m/s} عند
\ang{30}. (ج)~الهواء الذي يبدأ أبعد عن خطّ الاستواء يبدأ أقرب إلى المحور، بزخم زاوي فائض أقلّ: فيصل بريح غربية أضعف عند كل خطّ عرض. وتخفّض توزيع الفرع الصاعد \omterm{def:b1:atmospheric-circulation:jet}{التيّار النفّاث} المتوقّع، لكن ليس بما يكفي للوصول إلى \qty{40}{m/s} المرصودة.
\end{solution}

\begin{solution}{exo:b1:atmospheric-circulation:7}
(أ)~تهبّ الريح نحو الشمال الشرقي. ومن يقف ظهره إليها، يواجه الشمال الشرقي، فيكون المنخفض على يساره: نحو الشمال الغربي. (ب)~تهبّ الريح نحو الجنوب الشرقي؛ وفي نصف الكرة الجنوبي يكون المنخفض على يمين من يقف ظهره إلى الريح: نحو الجنوب الغربي. (ج)~تُحرف ريح السطح من الخطوط المتساوية الضغط نحو المنخفض، أي عكس عقارب الساعة من \omterm{def:b1:atmospheric-circulation:geostrophic}{الريح الجيوستروفية} في نصف الكرة الشمالي. \omterm{def:b1:atmospheric-circulation:geostrophic}{فالريح الجيوستروفية}، والمنخفض عمودي عليها إلى يسارها، يُداران إذن في عقارب الساعة عن إجابة (أ): يقع المنخفض بين الشمال الغربي والشمال، عند \omterm{def:b1:spherical-astronomy:horizontal}{السمت} $\ang{315} + \alpha$ لزاوية عبور
$\alpha$.
\end{solution}

\begin{solution}{exo:b1:atmospheric-circulation:8}
(أ)~المساحة بين خطّ الاستواء و\ang{30} شمالًا هي $2\pi a_\oplus^2\sin\ang{30} =
\qty{1.275e14}{m^2}$؛ وكتلة الهواء فوق متر مربّع بين المستويين هي
$\Delta p/g = \num{9e4}/9.81 = \qty{9174}{kg}$؛ والمجموع
\qty{1.170e18}{kg}. (ب)~$\num{1.170e18}/\num{8.7e10} = \qty{1.34e7}{s}$، أي نحو
\num{156} يومًا، خمسة أشهر. (ج)~$\num{1.170e18}/\num{21.4e10} = \qty{5.47e6}{s}$،
أي نحو \num{63} يومًا، شهرين.
\end{solution}

\begin{solution}{exo:b1:atmospheric-circulation:9}
(أ)~$\num{0.01} \times \num{86400} = \qty{864}{m}$ في اليوم. (ب)~$0.864 \times 9.8 =
\qty{8.5}{K}$ في اليوم. (ج)~يبرد الهواء بانبعاث الأشعة تحت الحمراء: يفقد الغلاف الجوي بالإشعاع طاقة أكثر مما يمتصّ، ويعوّض هبوط الهواء الخسارة بتسخينه بالضغط. (د)~الهواء الهابط يسخن ويجفّ بالنسبة إلى التشبّع، فتتبخّر سحبه ولا يمكن أن يتكوّن المطر: \omterm{def:b1:atmospheric-circulation:cells}{والمرتفعات شبه المدارية} أحزمة سماء صافية ومطر قليل، حيث تتشكّل الصحاري الحارّة على اليابسة.
\end{solution}

\begin{solution}{exo:b1:atmospheric-circulation:10}
(أ)~$V^2 + frV - fr|\vec v_g| = 0$ مع $fr = \qty{55.85}{m/s}$:
$V = \tfrac12\bigl(-55.85 + \sqrt{55.85^2 + 4 \times 55.85 \times 15}\bigr) =
\qty{12.3}{m/s}$، أي \qty{82}{\percent} من \omterm{def:b1:atmospheric-circulation:geostrophic}{الريح الجيوستروفية}. (ب)~$V^2 - frV +
fr|\vec v_g| = 0$ لها جذور حقيقية فقط إذا كان $(fr)^2 \geq 4fr|\vec v_g|$، أي
$|\vec v_g| \leq fr/4 = \qty{14.0}{m/s}$. (ج)~قرب مركز مرتفع، حيث $r$ صغيرة، يكون الحدّ $fr/4$ صغيرًا: يجب أن يكون تدرّج الضغط ضعيفًا، والخطوط المتساوية الضغط متباعدة، والرياح خفيفة. ولا يملك المنخفض مثل هذا الحدّ: فخطوطه المتساوية الضغط يمكن أن تتزاحم قرب المركز، مع رياح قوية.
\end{solution}

\begin{solution}{exo:b1:atmospheric-circulation:11}
(أ)~من أجل $\varphi$ صغيرة، $\sin\varphi \approx \varphi$ و$\cos\varphi \approx 1$، إذن $u \approx \Omega a_\oplus\varphi^2$. (ب)~$\ang{10} = \qty{0.1745}{rad}$:
$464.6 \times 0.1745^2 = \qty{14.15}{m/s}$، مقابل \qty{14.23}{m/s} بالضبط، أي \qty{0.5}{\percent} أقلّ. (ج)~بحلّ $464.6\sin^2\varphi/\cos\varphi = 27.0$ عدديًّا نجد $\varphi = \ang{13.7}$. (د)~تطابق الحلقة \omterm{def:b1:atmospheric-circulation:jet}{التيّار النفّاث} السنوي المتوسط الحقيقي عند أقلّ من \ang{14}: فحفظ الزخم الزاوي يصف الفرع العلوي \omterm{def:b1:atmospheric-circulation:cells}{لخليّة هادلي} في المدارات العميقة فقط؛ وشمال نحو \ang{15} تتولّى الدوّامات الأمر.
\end{solution}

\begin{solution}{exo:b1:atmospheric-circulation:12}
(أ)~المتوسط $31.7/12 = \ang{2.6}$ شمالًا؛ والمدى من \ang{12.8} جنوبًا إلى \ang{10.9} شمالًا، أي
\ang{23.7}. (ب)~يأتي الموضع الأقصى شمالًا في سبتمبر، بثلاثة أشهر بعد انقلاب يونيو؛ والأقصى جنوبًا في فبراير، بشهرين بعد انقلاب ديسمبر.
ويأتي التأخّر من تخزين الحرارة، وخاصة في \omterm{def:b1:fluid-envelopes-structure:ocean-layers}{الطبقة المختلطة} للمحيط (\cref{ch:b1:fluid-envelopes-structure})، التي تسخن وتبرد ببطء. (ج)~يتوسّط المتوسط الزونالي كل \omterm{def:b1:spherical-astronomy:geographic}{خطوط الطول}: فإذا امتدّ نطاق المطر فوق القارات بعيدًا إلى نصف الكرة الصيفي، فإن متوسط كل \omterm{def:b1:spherical-astronomy:geographic}{خطوط الطول} يمكن أن يعبر خطّ الاستواء مع أن النطاق فوق محيط واحد لا يعبره.
\end{solution}

\begin{solution}{pb:b1:atmospheric-circulation:1}
\begin{enumerate}
  \item $\rho = p/(R_dT) = \num{100600}/(287.05 \times 283.15) =
        \qty{1.238}{kg.m^{-3}}$.
  \item $|\nabla p| = \qty{400}{Pa}/\qty{2.5e5}{m} = \qty{1.6e-3}{Pa.m^{-1}}$.
  \item $|\nabla p|/\rho = \num{1.6e-3}/1.238 = \qty{1.293e-3}{m.s^{-2}}$.
  \item $\Omega = 2\pi/\num{86164.2} = \qty{7.292e-5}{rad.s^{-1}}$;
        $f = 2\Omega\sin\ang{50} = \qty{1.117e-4}{s^{-1}}$.
  \item $|\vec v_g| = \num{1.293e-3}/\num{1.117e-4} = \qty{11.57}{m/s}$.
  \item $11.57/0.514444 = \qty{22.5}{kn}$؛ $11.57 \times 3.6 = \qty{41.7}{km/h}$.
  \item الزمن هو $|\vec v_g|/(|\nabla p|/\rho)$، و$|\vec v_g| = (|\nabla
        p|/\rho)/f$، إذن هو $1/f = \qty{8951}{s}$، أي \qty{2.49}{h}؛ و$2\pi/f =
        \qty{15.6}{h}$. والأنماط التي تدوم يومًا أو أكثر قريبة إذن من التوازن.
  \item تشير \omterm{def:b1:atmospheric-circulation:pgf}{قوة تدرّج الضغط} نحو الشمال، نحو المنخفض. وتهبّ \omterm{def:b1:atmospheric-circulation:geostrophic}{الريح الجيوستروفية} على طول الخطوط المتساوية الضغط والمنخفض على يسارها: نحو الشرق، من الغرب، \omterm{def:b1:spherical-astronomy:horizontal}{السمت} \ang{270}: ريح غربية.
  \item إذا وقفت ظهرك إلى ريح غربية، مواجهًا الشرق، فإن يسارك هو الشمال، حيث يقع المنخفض: القانون محقَّق.
  \item العصا تشير غربًا من دائرة المحطة، نحو مصدر الريح؛ و\qty{22.5}{kn} تُقرَّب إلى \qty{20}{kn}، ريشتان كاملتان؛ وتقفان على الجانب مع عقارب الساعة من العصا، جانبها الشمالي، نحو المنخفض.
  \item عند \ang{50} جنوبًا، $f < 0$: تبقي الريح المنخفض على يمينها، فتهبّ نحو الغرب، من الشرق.
  \item $f(\ang{30}) = \qty{7.292e-5}{s^{-1}}$: \qty{17.7}{m/s}؛ $f(\ang{60}) =
        \qty{1.263e-4}{s^{-1}}$: \qty{10.2}{m/s}.
  \item $f(\ang{5}) = \qty{1.271e-5}{s^{-1}}$ يعطي \qty{101.7}{m/s}، ريحًا لم تُرصد قط بتدرّجات كهذه. زمن الدوران $1/f$ هو \qty{22}{h}: لا وقت للتوازن كي ينشأ، والصيغة لا تنطبق.
  \item $9.0/11.57 = 0.778$؛ على طول الخطوط المتساوية الضغط $9.0\cos\ang{25} = \qty{8.16}{m/s}$،
        وعبرها نحو المنخفض $9.0\sin\ang{25} = \qty{3.80}{m/s}$.
  \item تهبّ \omterm{def:b1:atmospheric-circulation:geostrophic}{الريح الجيوستروفية} نحو \omterm{def:b1:spherical-astronomy:horizontal}{السمت} \ang{90}؛ وحين تُدار بمقدار \ang{25} نحو المنخفض في الشمال، تهبّ ريح السطح نحو \ang{65}، من
        \ang{245}، أي من الغرب الجنوبي الغربي.
  \item $\tan\alpha = 1/(f\tau)$ يعطي $\tau = 1/(\num{1.117e-4}\tan\ang{25}) =
        \qty{1.92e4}{s}$، \qty{5.3}{h}؛ يتنبّأ النموذج ب$11.57\cos\ang{25} =
        \qty{10.5}{m/s}$، أكثر من \qty{9.0}{m/s} المرصودة: فالاحتكاك الحقيقي ليس خطّيًّا تمامًا في الريح، ويبطّئها أكثر من النموذج عند هذه الزاوية.
  \item عبر أسطوانة نصف قطرها $R = \qty{500}{km}$ وارتفاعها $h =
        \qty{1}{km}$: $\dot m = 2\pi R h\rho u_r = 2\pi \times \num{5e5} \times 1000
        \times 1.238 \times 3.80 = \qty{1.48e10}{kg/s}$.
  \item $w = \dot m/(\rho\pi R^2) = 2hu_r/R = 2 \times 1000 \times 3.80/\num{5e5} =
        \qty{0.0152}{m/s}$، نحو \qty{1310}{m} في اليوم.
  \item $1.31 \times 9.8 = \qty{12.9}{K}$ في اليوم: يبلغ الهواء التشبّع قريبًا، فيمتلئ المنخفض بالسحب والمطر. وتحت مرتفع جوّي يتدفّق هواء السطح نحو الخارج، ويهبط الهواء الذي فوقه ويسخن، وتصحو السماء.
  \item $\Omega a_\oplus = \num{7.292e-5} \times \num{6.371e6} = \qty{464.6}{m/s}$
        عند خطّ الاستواء؛ $464.6\cos\ang{30} = \qty{402.3}{m/s}$ عند \ang{30} شمالًا.
  \item $u = 464.6\,\sin^2\varphi/\cos\varphi$: \qty{57.8}{m/s} عند \ang{20} شمالًا و
        \qty{134.1}{m/s} عند \ang{30} شمالًا.
  \item $u = 464.6\,(\cos^2\ang{10} - \cos^2\ang{30})/\cos\ang{30} =
        \qty{117.9}{m/s}$.
  \item $134.1/42.3 = 3.17$. تأخذ دوّامات \omterm{def:b1:spherical-astronomy:geographic}{خطوط العرض} الوسطى، أي المنخفضات والمرتفعات الجوّية، الزخم الزاوي من الفرع العلوي وتحمله إلى مكان آخر؛ ويصعد هواء الخليّة فوق نطاق من \omterm{def:b1:spherical-astronomy:geographic}{خطوط العرض} بدلًا من أن يكون ساكنًا عند خطّ الاستواء. ويتوسّط المتوسط الزونالي أيضًا تيّارًا متعرّجًا فوق كل \omterm{def:b1:spherical-astronomy:geographic}{خطوط الطول}.
  \item $M/M_0 = (402.3 + 42.3)\cos\ang{30}/464.6 = 0.829$: فقد الهواء
        \qty{17}{\percent} من الزخم الزاوي الذي كان سيحتفظ به لو بقي حلقة حرّة.
  \item \omterm{def:b1:atmospheric-circulation:geostrophic}{الريح الجيوستروفية} عند P هي \qty{11.6}{m/s}، من الغرب؛ والريح المقاسة عند الأرض \qty{9.0}{m/s}، أي نحو أربعة أخماسها، ومُحرّفة
        \ang{25} نحو المنخفض بالاحتكاك.
\end{enumerate}
\end{solution}
